\textbf{T3: Dry ice hockey (10 pts)}

At pressure $p_0 = 100\,\mathrm{kPa}$, solid CO$_2$ (dry ice) sublimates (goes
from solid to gaseous state) at $T_s = -78.5\,{}^\circ\mathrm{C}$. Its
saturated vapour pressure follows the Clausius–Clapeyron relation:
\[
\frac{\mathrm{d}p_\mathrm{sat}}{\mathrm{d}T} = \frac{\mu\lambda p_\mathrm{sat}}{RT^2},
\]
where the latent heat of sublimation is $\lambda = 600\,\mathrm{kJ/kg}$, the
molar mass is $\mu = 0.044\,\mathrm{kg/mol}$, and the gas constant is
$R = 8.3\,\mathrm{J/(mol \cdot K)}$. The thermal conductivity of the CO$_2$
gas is $\kappa = 10\,\mathrm{mW/(m \cdot K)}$ and its dynamic viscosity is
$\eta = 10\,\upmu\mathrm{Pa \cdot s}$. The density of dry ice is
$\rho = 1500\,\mathrm{kg/m^3}$ and the gravitational acceleration is
$g = 10\,\mathrm{m/s^2}$.

\begin{center}\includegraphics[width=0.44\linewidth]{figures/EuPhO_2026_Q3_fig1.png}\end{center}

A puck of radius $r = 10\,\mathrm{mm}$ consists of a disc of dry ice of
thickness $h = 1\,\mathrm{mm}$, and a metal disc of mass $M = 0.01\,\mathrm{kg}$
on top of it. The initial temperature of the puck is $T_s$ and the ambient
pressure is $p_0$.

The puck is placed on a horizontal metal plate which is held at a constant
temperature $T = T_s + \Delta T$, and given an initial horizontal velocity
$v = 10\,\mathrm{mm/s}$. After a very long time, the horizontal displacement
$L$ of the metal disc is measured.

a) \textbf{(2 pts)} When $\Delta T$ is sufficiently small, $L$ is negligible
and independent of $\Delta T$. However, when $\Delta T$ reaches a critical
value $\Delta T_c$, the function $L(\Delta T)$ starts to grow. Estimate
$\Delta T_c$.

b) \textbf{(8 pts)} Estimate the maximal value $L_\mathrm{max}$ of the function
$L(\Delta T)$.

Treat the CO$_2$ gas as ideal. Assume no tilting of the puck at any moment, all
surfaces are perfectly smooth, and the thermal conductivities of the metal, dry
ice, and gas satisfy $\kappa_\mathrm{metal} \gg \kappa_\mathrm{ice} \gg \kappa$.
